\documentclass[10pt]{letter}
\usepackage[a4paper,margin=1.7cm,top=1.4cm]{geometry}
\usepackage[T1]{fontenc}
\usepackage{lmodern}
\signature{Mohd-Zulhilmi Paiz Ismadi\\Department of Mechanical Engineering\\School of Engineering, Monash University Malaysia\\47500 Bandar Sunway, Selangor, Malaysia\\mohd.zulhilmi@monash.edu}
\address{}
\date{\today}
\setlength{\longindentation}{0pt}
\setlength{\indentedwidth}{\textwidth}
\begin{document}
\begin{letter}{Guest Editors, Special Issue ``Integrating Imaging with Personalized Cardiovascular Medicine: Current
Advances and Future Directions''\\IEEE Journal of Biomedical and Health Informatics}
\opening{Dear Guest Editors,}

We submit the manuscript ``Topological Segmentation Error and Boundary-Condition Tuning in Computed Coronary FFR: A
Controlled In Silico Study'' for consideration in the special issue. The call identifies uncertainty propagation,
segmentation errors and robust validation frameworks as central challenges for image-based cardiovascular models. Our
study addresses all three, for computed fractional flow reserve (FFR).

Computed FFR depends on the segmented coronary lumen twice, through the arteries in the model and through outlet boundary conditions that are derived from those arteries and, increasingly, tuned to match perfusion imaging. We
inserted 150 controlled stenoses into 108 coronary trees from a public CT angiography dataset, applied four segmentation
error types and recomputed FFR under fixed, re-derived and perfusion-tuned boundary conditions (one global or one
per-territory tuning parameter) in two microvascular bed structures, with a three-dimensional case study and a
replication at doubled and tripled hyperemic demand.

\textbf{Innovation.} The study is new in three respects. First, it judges segmentation error by whether it changes
the treatment decision at FFR 0.80, against the variability of repeat invasive FFR, rather than by a continuous
sensitivity. Second, it solves the same corrupted anatomy under fixed, re-derived and tuned boundary conditions,
which isolates the effect of tuning, and it separates topological errors (missing or broken vessels) from caliber
errors of measured inter-observer size. Third, it quantifies, to our knowledge for the first time, how often a model
passes a perfusion check while its FFR is materially wrong, and tests the reduced-order result with
three-dimensional computational fluid dynamics.

\textbf{Main findings.}
\begin{itemize}
\item Topological errors (a missed side branch or a vessel break) changed the treatment decision at 0.80 in 32--33\%
of models with fixed boundary conditions, compared with 6--10\% for caliber errors (a vessel drawn too wide or too narrow) of inter-observer magnitude, which
were close to the variability of repeat invasive measurement.
\item On identical anatomy, a vessel break changed the decision in 43\% of models with fixed boundary conditions and
in 5\% with re-derived ones.
\item Tuning the boundary conditions to perfusion reduced decision changes, and per-territory tuning corrected the
typical missed branch, as did the three-dimensional case. A tuned model, however, passes the perfusion check whether
or not its FFR is correct: 5--20\% of tuned topological-error models and 4--18\% of tuned caliber-error models passed
a check stricter than measurement repeatability while their FFR was wrong by more than 0.05, against 3--4\% for
models with correct anatomy.
\end{itemize}

\textbf{Significance.} These results bear directly on how image-based models are validated. Agreement with
perfusion after tuning does not establish that the computed FFR is correct, and overlap-based segmentation metrics do not capture the error types that
carry the decision risk. The work combines medical image analysis, physiological modeling and decision-level
evaluation, which matches the scope of the Journal and of this special issue. 

The manuscript is original, has not been published, and is not under consideration elsewhere. All authors have read
and approved the submission. The authors declare no conflicts of interest, and the work received no external funding.
The study used a public, anonymized dataset (ImageCAS-X, CC BY 4.0). The code is available from the authors on request. Supplementary Material accompanies the manuscript.

\closing{Yours sincerely,}
\end{letter}
\end{document}
